\chapter{Dữ liệu}

\section{Nguồn dữ liệu và độ tin cậy}
\begin{longtable}{p{0.18\textwidth}p{0.26\textwidth}p{0.16\textwidth}p{0.28\textwidth}}
\toprule
\textbf{Nguồn} & \textbf{Đơn vị cung cấp/URL} & \textbf{Ngày truy cập} & \textbf{Lý do tin cậy và giấy phép} \\
\midrule
\textit{Nguồn 1} & \textit{Đơn vị và liên kết} & \textit{Ngày} & \textit{Mô tả} \\
\textit{Nguồn 2} & \textit{Đơn vị và liên kết} & \textit{Ngày} & \textit{Mô tả} \\
\bottomrule
\end{longtable}

\section{Mô tả bộ dữ liệu}
Nêu đơn vị quan sát, số dòng, số cột, phạm vi thời gian, phạm vi địa lý và định dạng.
Xác nhận bộ dữ liệu có tối thiểu 2.000 dòng và 7 biến độc lập.

\section{Data dictionary}
Trình bày các biến quan trọng tại \cref{app:data-dictionary}; phân biệt biến gốc,
biến dẫn xuất, biến định lượng và định tính.

\section{Thu thập và tái tạo dữ liệu}
Mô tả cách tải dữ liệu, tham số, phiên bản, cấu trúc thư mục và lệnh chạy pipeline.
Nêu cách bảo đảm người khác có thể tái tạo bộ dữ liệu phân tích.

\section{Làm sạch và biến đổi}
\begin{itemize}
    \item Chuẩn hóa kiểu dữ liệu, tên biến, đơn vị và địa danh Việt Nam.
    \item Xử lý giá trị thiếu, ngoại lệ và bản ghi trùng lặp.
    \item Kết hợp các nguồn và tạo biến dẫn xuất.
    \item Bảo toàn dữ liệu gốc và lưu dấu vết xử lý.
\end{itemize}

\section{Đánh giá chất lượng dữ liệu}
Trình bày tỷ lệ thiếu, trùng lặp, ngoại lệ, kiểm tra miền giá trị và các giới hạn.
Giải thích tác động của chất lượng dữ liệu lên kết luận.

\section{Đạo đức, quyền riêng tư và hạn chế}
Nêu giấy phép, quyền riêng tư, thiên lệch dữ liệu và các trường hợp không nên diễn
giải quá mức.
